```latex
\documentclass{article}
\usepackage{pathfinder-note}

\title{A CVaR Oracle for a Balanced Allocation Game}

\begin{document}
\maketitle

\section{The pair}

Q designs quadratic payments for a resource allocation game with separate agent decisions and a shared capacity constraint. P studies sampled conditional value at risk (CVaR) optimization over a communication network in which every agent maintains a copy of the same decision. \ledger{1, 3}

\section{Connection pursued}

The question was whether P's black-box CVaR estimates could supply risk-adjusted valuations for Q's incentive game, with guarantees for the true CVaR objective. P's consensus iteration cannot be transferred directly: consensus is appropriate for its shared decision, whereas Q's welfare optimum can require unequal allocations. P's assumption that the feasible set contains a ball centered at zero also cannot be copied literally to Q's nonnegative allocation sets. The reusable ingredient is the local empirical-CVaR estimate. \ledger{3, 20}

\section{What was found}

\begin{objection}
Q's stated strict uniqueness certificate is incompatible with its implementation conditions. P2(i) requires \(\sum_m A_{nm}^n=0\). In the direction with zero allocation change and the same nonzero price change for every agent, this makes the quadratic form of Q's pseudogradient bound zero, contradicting P1(ii)'s strict negative bound. Thus Q's proposed payment cannot satisfy both stated LMIs. This does not, by itself, rule out implementation. \ledger{5, 6}
\end{objection}

For Q's payment in Eq.~(40), implementation has a direct argument. Write \(P^k=\sum_n p_n^k\) and \(z^k=\sum_n x_n^k-c^k\). Its price best response satisfies
\[
\frac{\partial U_n}{\partial p_n^k}
=-\frac{\alpha}{N-1}(Np_n^k-P^k-z^k).
\]
If one price in coordinate \(k\) is positive, every price there must be positive and equal, and summing the interior conditions gives \(z^k=0\). If all prices vanish, the boundary conditions give \(z^k\leq0\). Hence every equilibrium has a common price \(p\), feasible allocations, and \(p^\top(\sum_nx_n-c)=0\). At that price, each allocation best response maximizes \(V_n(x_n)-\alpha p^\top x_n\), exactly the planner's condition with multiplier \(\lambda=\alpha p\). Under Q's strong concavity and existence of a planner multiplier, the converse follows from joint concavity of each player's payoff. The allocation is unique, but the multiplier and equilibrium price need not be. \ledger{8, 9, 14}

For example, with two agents, one resource, \(c=2\), \(\mathcal X_n=[0,1]\), \(\alpha=1\), and \(V_n(x)=2x-x^2/2\), the allocation \((1,1)\) is an equilibrium with any common price in \([0,1]\). Q's unique-message claim therefore fails. Its original payment can also violate individual rationality: with \(c=1\), \(V_1(x)=4x-x^2/2\), and \(V_2(x)=x-x^2/2\), the equilibrium \((x_1,x_2,p_1,p_2)=(1,0,3,3)\) charges agent 1 \(9/2\), exceeding \(V_1(1)=7/2\). Both examples can be realized by bounded, convex stochastic losses of the type allowed in P. \ledger{7, 10}

The ledger gives an others-only payment correction. Let
\[
\kappa_N=1+\frac{N}{(N-1)^2},
\qquad
g_n(s_{-n})=
\frac{\alpha N}{(N-1)^3}
\bar p_{-n}^{\top}
\left(\bar x_{-n}-\frac{N-1}{N}c\right),
\]
and replace Q's payment \(t_n\) by \(t_n+g_n\). Since \(g_n\) contains no message from agent \(n\), best responses and equilibria are unchanged. Substitution at a common-price, complementary equilibrium gives
\[
t_n=\kappa_N\lambda^\top(x_n-c/N),
\qquad
t_n+g_n=\lambda^\top(x_n-c/N).
\]
The corrected payments sum to zero. If \(0\in\mathcal X_n\) and \(V_n(0)=0\), the allocation best response also gives
\[
V_n(x_n)-\lambda^\top x_n+\lambda^\top c/N\geq0,
\]
so the corrected mechanism is individually rational at every equilibrium. \ledger{12, 13, 14, 24}

For the CVaR connection, let \(C_n(x)\) be agent \(n\)'s true CVaR loss and \(H_n(x)\) its expected empirical CVaR after ball smoothing with radius \(\delta_n\) and a fresh batch of size \(s_n\). P's bounded-loss CVaR comparison and integrated empirical-distribution bound give the pointwise estimate
\[
|H_n(x)-C_n(x)|\leq h_n,
\qquad
h_n=L_n\delta_n+
\frac{U_n\sqrt{2\pi}}{\beta_n\sqrt{s_n}},
\]
provided the perturbation queries are available on the required neighborhood. Form true and surrogate valuations by subtracting \(C_n\) and \(H_n\), respectively, from the same strongly concave utility term, then normalize each valuation at zero. The surrogate valuations retain the required strong concavity. \ledger{15, 17, 24, 26}

At any equilibrium of the corrected surrogate game, the true planner welfare loss is at most \(2\sum_nh_n\), true unilateral regret is at most \(2h_n\) for agent \(n\), and true utility is at least \(-2h_n\) relative to its normalized outside option. Payments balance exactly. The welfare and regret bounds compare two decisions, so their normalization constants cancel; they follow from the two-sided CVaR value error and surrogate optimality. Strong concavity also gives an allocation-distance bound \(2\sqrt{\sum_nh_n/\alpha}\). These are static, conditional guarantees. \ledger{17, 18, 22, 24}

\section{Objections and limits}

The payment correction resolves the original individual-rationality counterexample without changing strategic incentives. It does not restore Q's strict LMI or unique prices. P's convergence theorem concerns consensus on one shared decision and does not prove convergence of sampled play to Q's equilibrium set. Q's numerical sample rule, \(Q_i=\lceil0.96^{i+1}\rceil\), equals one at every iteration as written, whereas its convergence theorem requires geometrically increasing sample sizes. Q also guarantees capacity feasibility only at equilibrium. Feasible perturbation queries near nonnegative boundaries, multiplier existence, and an implementable learning rule remain conditions or open tasks. \ledger{19, 20, 22}

\section{Outside sources and next step}

A narrow prior-work search reported adjacent CVaR market work at \url{https://arxiv.org/abs/2003.06119} and stochastic aggregative-game learning at \url{https://arxiv.org/abs/2501.06023}. Neither source was used to prove the claims above, and the search does not establish novelty. \ledger{16, 20}

The smallest decisive next step is to prove or refute convergence of a feasible, sample-based learning rule to the corrected game's equilibrium set, beginning with two agents and one resource. The present material establishes a conditional static connection; it is thin on dynamics and cannot support a publication-level novelty claim. \ledger{20, 22, 24}

\end{document}
```